\section{The original escaping endpoint and the arithmetic heat family}
\label{sec:eh}
\newcounter{eh}
\renewcommand{\theeh}{EH\arabic{eh}}
\newcommand{\EH}{\refstepcounter{eh}\par\medskip\noindent\textbf{\theeh.}\ }

\EH\textbf{Sources and exact scope.}
The two supplied versions of \emph{The flip, the Fable fibre on the Riemann
sphere, and Weil positivity} have identical SHA256 values; they are retained
as \path{inputs/flip_a.txt} and \path{inputs/flip_b.txt}. Their supplied
attribution is Claude (Anthropic), dated 23 September 2026. They are inputs
to this calculation, not primary literature certificates. The original
programme calculation of the escaping orbit is retained in
\path{material_generator_endpoint.md} and
\path{tex/inverse_fibre_heat.tex} in the 8 September fluid research edition.
All identities used from that orbit are proved again below. Its time is
written $\tau_{\rm f}$ here to distinguish it from the adjoined semiring
zero $\tau$. This is an explicit renaming, not a change of its time scale.
The original spatial third coordinate, called $t$ in the polynomial, is
written $w$. No coefficient is changed.

Rodgers--Tao's original author source defines
\begin{align}
 H_t(Z)&=\int_0^\infty e^{tu^2}\Phi(u)\cos(Zu)\,du,
 &H_0(Z)&=\frac18\xi\left(\frac12+\frac{iZ}{2}\right),\label{eq:eh-heat}\\
 \Phi(u)&=\sum_{n\ge1}(2\pi^2n^4e^{9u}-3\pi n^2e^{5u})
 e^{-\pi n^2e^{4u}},
 &\xi(s)&=\frac{s(s-1)}2\pi^{-s/2}\Gamma(s/2)\zeta(s).\notag
\end{align}
These are equations \texttt{hoz}, \texttt{sas}, \texttt{phidef},
\texttt{htdef} in \href{https://arxiv.org/abs/1801.05914v5}{Rodgers--Tao,
The de Bruijn--Newman constant is non-negative}, original TeX lines78--103.
The factor $1/8$ in this definition is a nonzero scalar; the number $1/8$
in the escaping trajectory below is a time. Their different roles are
resolved by an explicit map, not by matching the displayed numerals.

\EH\textbf{The original polynomial and its inverse-root chart.}
Let $F:\mathbb C^3_{x,y,w}\to\mathbb C^3_{a,b,c}$ be
\begin{align*}
 F_1&=(1+xy)^3w+y^2(1+xy)(4+3xy),\\
 F_2&=y+3x(1+xy)^2w+3xy^2(4+3xy),\\
 F_3&=2x-3x^2y-x^3w.
\end{align*}
On $x\ne0$ set $v=1+xy$, $d=2-3xy-x^2w$, and
$A=v^2+v-v^3d$, $B=4v+2-3v^2d$. Then
$F=(A/x^2,B/x,xd)$ and the coordinate Jacobian
$\partial(x,v,d)/\partial(x,y,w)$ has determinant $-x^3$.
The determinant of $\partial F/\partial(x,v,d)$ is $x^{-3}$ times
\[
\det\begin{pmatrix}-2A&2v+1-3v^2d&-v^3\\
-B&4-6vd&-3v^2\\d&0&1\end{pmatrix}=2.
\]
Expansion reduces the last expression to
$(4-6vd)(-2v^2-2v+3v^3d)+(2v+1-3v^2d)(4v+2-6v^2d)=2$.
Hence $\det DF=-2$ on a dense open set and therefore everywhere.

Put $P_{a,b,c}(r)=cr^3-2r^2+br-2a$. For $x\ne0$, the exact inverse
coordinates are
\[
 r=y+x^{-1},\quad \alpha=x^{-1},\quad
 (x,y,w)=(\alpha^{-1},r-\alpha,5\alpha^2-3r\alpha-c\alpha^3).
\]
Substitution gives
\[
F=(r^2+r\alpha-cr^3,\ 4r+2\alpha-3cr^2,\ c).
\]
Thus $F=(a,b,c)$ is equivalent to $P(r)=0$ and
$\alpha=P'(r)/2\ne0$. This proves both directions, including the
excluded multiple-root locus. On $x=0$ the formula is
$F(0,y,w)=(w+4y^2,y,0)$, so the entire additional inverse point is
$(0,b,a-4b^2)$ when $c=0$.

\EH\textbf{The escaping trajectory, with its original endpoint.}
For $\tau_{\rm f}<1/8$ put
\[
 z=\sqrt{1-8\tau_{\rm f}}>0,\qquad
 \gamma(\tau_{\rm f})=(z^{-1},-3z/2,13z^2/2).
\]
Direct substitution in the original polynomial gives
\[
 F(\gamma)=(-z^2/4,0,0)=(-1/4+2\tau_{\rm f},0,0),\quad
 r(\gamma)=-z/2,\quad \alpha(\gamma)=z.
\]
Moreover $z'=-4/z$. Differentiating the displayed identity gives
$DF(\gamma)\gamma'=(2,0,0)^{\mathsf T}$. For
\[
 U=2(DF)^{-1}e_1+6F_3(DF)^{-1}e_2
\]
this proves $\gamma'=U(\gamma)$, since $F_3(\gamma)=0$ and $DF$ is
invertible. The trajectory has no finite affine limit because its first
coordinate tends to infinity. Its target has the finite limit $(0,0,0)$.
At that target the entire affine fibre is $(0,0,0)$, by the preceding
inverse description: the double root $r=0$ has $\alpha=0$ and is excluded
from the $x\ne0$ chart. These are exact statements about this field;
they require no claim concerning a separate Navier--Stokes solution.

\EH\textbf{The exact heat map and its time origin.}
Retain the incoming note's arithmetic coordinate
$s=1/2+r$ and heat coordinate $Z=-2ir$. Define
\[
 h=4\tau_{\rm f}-\frac12=2a,\qquad
 Q_h(Z)=Z^2-2h.
\]
On $b=c=0$ one has the polynomial identity
\[
 Q_h(Z)=2P_{a,0,0}(r)\quad (Z=-2ir,\ h=2a).
\]
Consequently $\partial_hQ_h=-2=-\partial_Z^2Q_h$, with exactly the
backward heat sign of \eqref{eq:eh-heat}. Along the escaping branch
$Z=iz$ and $h=-z^2/2$. The original times map as
\[
 \tau_{\rm f}=0\longmapsto h=-1/2,\qquad
 \tau_{\rm f}\uparrow1/8\longmapsto h\uparrow0.
\]
There is no step identifying $h$ with the absolute arithmetic time $t$
of \eqref{eq:eh-heat}; the formula specifies the translated and scaled
time of this particular polynomial solution completely.
Its zeros are $\pm i\sqrt{-2h}$ for $h<0$, a double zero at $0$ for
$h=0$, and $\pm\sqrt{2h}$ for $h>0$. Thus its all-real threshold is
exactly $h=0$, including the endpoint. The escape reaches that real
double root. It does not leave a nonreal root at the endpoint.

The reciprocal of the escaping affine coordinate has an exact relation
to the meromorphic heat-root velocity. On the same $b=c=0$ chart,
$P'(r)=-4r$, so the inverse formulas give
\begin{equation}
 \alpha=-2r=-iZ,\qquad x=\frac{i}{Z},\qquad
 \frac{dZ}{dh}=\frac1Z,\qquad
 x=i\frac{dZ}{dh}.
 \label{eq:eh-reciprocal-bridge}
\end{equation}
Here the derivatives are along either root branch on $Z\ne0$:
differentiating $Z^2=2h$ proves $2Z\,dZ/dh=2$. The original root
coordinate and original flow time therefore satisfy, with no change
of scale left implicit,
\begin{equation}
 h=-2r^2,\qquad \frac{dr}{dh}=-\frac1{4r},\qquad
 \frac{dr}{d\tau_{\rm f}}=-\frac1r,
 \qquad \frac{dh}{d\tau_{\rm f}}=4.
 \label{eq:eh-original-reciprocal}
\end{equation}
The ordered-pair rational energy of the two heat roots $Z,-Z$ is
\begin{equation}
 \mathcal E=\frac{2}{(2Z)^2}
 =\frac1{2Z^2}=-\frac1{8r^2}=-\frac{x^2}{2}.
 \label{eq:eh-energy-coordinates}
\end{equation}
Each equality follows from $Z=-2ir$ and $x=i/Z$. In particular, on
the original escaping branch this rational function is
\begin{equation}
 \mathcal E(\gamma(\tau_{\rm f}))
 =-\frac1{2z^2}=-\frac1{2(1-8\tau_{\rm f})}
 \longrightarrow-\infty.
 \label{eq:eh-incoming-energy}
\end{equation}
This is its analytic continuation along imaginary heat roots $Z=iz$.
On the real-root branch $h>0$ the same rational function equals
$1/(4h)>0$ instead. Thus the reciprocal bridge retains both the
factor $4$ in time and the negative sign on the incoming branch.
These are calculations on $Q_h$; the arithmetic time $t$ in
\eqref{eq:eh-heat} has not been identified with $h$.

\EH\textbf{Retaining the collision algebra and its holonomy.}
The family
\[
 \mathcal A=\mathbb C[h,Z]/(Z^2-2h)
\]
is free of rank two over $\mathbb C[h]$, with basis $(1,Z)$: division
by the monic polynomial gives existence and uniqueness of this basis.
Its zero fibre is $\mathbb C[\eta]/(\eta^2)$ by $Z\mapsto\eta$.
For $h\ne0$ its two roots are exchanged by one loop around $h=0$,
since $\sqrt{2h_0e^{i\theta}}=\sqrt{2h_0}e^{i\theta/2}$ for
$0\le\theta\le2\pi$. On the two labelled roots this is the matrix
\[
 M_{\rm loop}=\begin{pmatrix}0&1\\1&0\end{pmatrix},\qquad
 M_{\rm loop}^*M_{\rm loop}=I,\quad M_{\rm loop}\ne I.
\]
This computes nontrivial \emph{unitary} monodromy around a fibre whose
only root is real. It supplies an explicit example against the implication
``any nonidentity holonomy forces a nonreal zero''. The map from the
family to this holonomy representation is the actual root covering over
$\mathbb C^*$, not an asserted arithmetic spectral identification.

\EH\textbf{The form through the collision.}
For real $h$, conjugation on $\mathcal A_h$ fixes $Z$. The Hermitian
trace form $q_h(f,g)=\operatorname{Tr}_{\mathcal A_h/\mathbb C}
(f\overline g)$ has matrix $\operatorname{diag}(2,4h)$ in $(1,Z)$.
Indeed multiplication by $Z$ has matrix
$\left(\begin{smallmatrix}0&2h\\1&0\end{smallmatrix}\right)$,
so its trace is zero and the trace of $Z^2$ is $4h$.
The signature is $(2,0)$ for $h>0$, $(1,1)$ for $h<0$, and a positive
line plus the nilpotent radical at $h=0$. In the original coordinate
$r=iZ/2$ with involution $r\mapsto-r$, its matrix is
$\operatorname{diag}(2,h)$: multiplying the second basis vector by
$i/2$ gives precisely that congruence. On point values the reflection
$s\mapsto1-\overline s$ sends $Z$ to $\overline Z$, because
$s=1/2+iZ/2$. Its induced conjugate-linear coordinate-algebra
involution, given by $f^\dagger(s)=\overline{f(1-\overline s)}$,
fixes the formal generator $Z=-2i(s-1/2)$:
\[
 Z^\dagger=2i(1/2-s)=-2i(s-1/2)=Z,
 \qquad r^\dagger=-r.
\]
Every sign and the degeneracy at zero therefore agree in the two
presentations. No strict positivity is asserted on the nilpotent radical.

\EH\textbf{The retained split boundary of the escaping chart.}
Let $L$ be the original bounded distributive support lattice and let
$G_L$ denote the original split-support semiring. Let
\[
 B=G_L(\mathbb C[z]),\quad B'=G_L(\mathbb C[z,z^{-1}]),\quad
 C=G_L(\mathbb C),
\]
where $B\to C$ is evaluation $z\mapsto e$ (supported amplitude zero).
The chart map to the trajectory is the ring homomorphism
\[
 \mathbb C[x,y,w]\longrightarrow\mathbb C[z,z^{-1}],\qquad
 (x,y,w)\longmapsto(z^{-1},-3z/2,13z^2/2),
\]
and applying $G_L$ gives its exact support-preserving semiring map.
Since $B'$ is the localization of $B$ at the supported polynomial $z$,
base change along $z\mapsto e$ proves
\begin{equation}
 B'\amalg_B C\ \cong\ C[e^{-1}]\ \cong\ L.\label{eq:eh-boundary}
\end{equation}
For clarity, the second isomorphism follows directly: $e^2=e$ makes
the image of $e$ equal to the unit when $e$ is invertible. Then
$ea=e$ for every supported amplitude $a$ makes all those amplitudes
equal to the top support, while each lower support label stays itself.
The support character $\chi_L$ witnesses that no two distinct labels
are identified. These equalities also prove the universal property.
The construction has thus retained a whole support spectrum at the
escaping endpoint. It has not created an affine complex value for
$z^{-1}$ at $z=0$: any attempted nonzero ring-valued base change
would send the invertible supported $e$ to ring zero and force $1=0$.
The exact receiver is $L$, with all its labels, as in
\eqref{eq:eh-boundary}.

\EH\textbf{The actual arithmetic coefficients, without a finite-packet replacement.}
For every $u\ge0$ each summand of $\Phi$ is strictly positive, since
\[
 2\pi^2n^4e^{9u}-3\pi n^2e^{5u}
 =\pi n^2e^{5u}(2\pi n^2e^{4u}-3)>0.
\]
For $a=\pi e^{4u}\ge\pi$,
$\sum_{n\ge1}n^4 e^{-an^2}\le
e^{-a/2}\sum_{n\ge1}n^4e^{-\pi n^2/2}$.
Hence $0<\Phi(u)\le C e^{9u}e^{-\pi e^{4u}/2}$ with a finite
constant $C$. This bound is integrable after multiplication by
any $u^N e^{T u^2+Ru}$, for any fixed $N,T,R\ge0$.
Uniform domination on compact subsets of $(t,Z)\in\mathbb C^2$
therefore proves that \eqref{eq:eh-heat} is entire jointly in $(t,Z)$,
and permits all the following differentiations and series integrations.

Define the exact positive moments
$M_j=\int_0^\infty u^{2j}\Phi(u)\,du>0$ for $j\ge0$. The full
Taylor expansion, with every time and spatial coefficient retained, is
\begin{equation}
 H_t(Z)=\sum_{m,n\ge0}
 \frac{(-1)^n M_{m+n}}{m!(2n)!}\,t^m Z^{2n}.\label{eq:eh-moments}
\end{equation}
Indeed the product of the absolute exponential and cosine Taylor series
is bounded on $|t|\le T,|Z|\le R$ by $e^{Tu^2+Ru}$, so Tonelli and
dominated convergence prove the equality. All displayed coefficients
are nonzero. All coefficients of odd powers of $Z$ are zero.
The identity $\partial_tH=-\partial_Z^2H$ follows coefficient by
coefficient, including the factorials, or directly under the integral.

The actual theta reciprocal can now be compared to
\eqref{eq:eh-reciprocal-bridge}, including its first discrepancy.
Fix a real arithmetic time $t$ and define
$M_j(t)=\int_0^\infty u^{2j}e^{tu^2}\Phi(u)\,du$. The same
domination proves that each moment is finite, and strict positivity
of the integrand on $u>0$ gives $M_j(t)>0$. Taylor expansion at
the spatial origin gives
\begin{equation}
 H_t(Z)=M_0(t)-\frac{M_1(t)}2Z^2
          +\frac{M_2(t)}{24}Z^4+O(Z^6).
 \label{eq:eh-theta-origin-jet}
\end{equation}
Consequently $H_Z/Z$ is holomorphic and nonzero at $Z=0$, with
value $-M_1(t)$. It is a holomorphic unit in a neighbourhood of
zero. Division in that punctured neighbourhood yields the exact
identity and Laurent expansion
\begin{align}
 \frac{H_{ZZ}}{H_Z}-\frac1Z
 &=\frac{ZH_{ZZ}-H_Z}{ZH_Z}
   =-\frac{M_2(t)}{3M_1(t)}Z+O(Z^3),
 \label{eq:eh-theta-reciprocal-defect}\\
 \frac{H_{ZZ}}{H_Z}
 &=\frac1Z-\frac{M_2(t)}{3M_1(t)}Z+O(Z^3).
 \label{eq:eh-theta-reciprocal-expansion}
\end{align}
For completeness, differentiating \eqref{eq:eh-theta-origin-jet}
gives $H_Z=-M_1(t)Z+M_2(t)Z^3/6+O(Z^5)$ and
$H_{ZZ}=-M_1(t)+M_2(t)Z^2/2+O(Z^4)$. Subtracting cancels the
linear numerator terms and gives
$ZH_{ZZ}-H_Z=M_2(t)Z^3/3+O(Z^5)$; its denominator is
$ZH_Z=-M_1(t)Z^2+O(Z^4)$. This proves the coefficient and its
negative sign in \eqref{eq:eh-theta-reciprocal-defect}.

For the quadratic, $(Q_h)_{ZZ}/(Q_h)_Z=1/Z$ identically wherever
$Z\ne0$. The theta correction just computed is nonzero because
$M_2(t)/M_1(t)>0$. Both quotients have a simple pole of residue
one at the spatial origin, but $H_t(0)=M_0(t)>0$ for every real
$t$, including $t=0$. Thus this pole is at a stationary point of
the actual theta function, and is not a zero of that function.
The quotient is a zero velocity only after restricting to an actual
simple zero. The displayed identities compute the exact discrepancy
between the two meromorphic quotients without identifying their
zero spaces or their absolute time coordinates.

\EH\textbf{The exact coefficient-support defect for this particular source.}
Represent $H_0$ initially as a series independent of $t$, with top
support on every even $Z$ coefficient in time degree zero, and bottom
support elsewhere. In the coefficient-labelled semiring
$(G_L(\mathbb C))[[t,Z]]$, lift a nonzero scalar coefficient to its
top-supported value and let the heat operators retain addition's join
of supports. The minus heat operator produces the coefficients in
\eqref{eq:eh-moments}. Its support is top at every $(m,2n)$ and bottom
at every $(m,2n+1)$.

Applying the plus heat operator to this series cancels the amplitude of
every positive time degree. At $(m,2n)$ that amplitude is
\[
 \frac{(-1)^nM_{m+n}}{(2n)!}
 \sum_{j=0}^m\frac{(-1)^j}{j!(m-j)!},
\]
which is zero for $m>0$ by $(1-1)^m=0$, and equals the original
$(-1)^nM_n/(2n)!$ for $m=0$. For $m>0$ every summand has top
support, so the cancelled coefficient is the supported zero $e$, not
absence $\tau$. All odd coefficients remain absent. This is a
calculated support memory for the actual theta source. Evaluation at
$t=0$ retains exactly the original $H_0$ coefficients. The general
closure and its inverse-on-fixed-support interpretation are proved
in the accompanying semilattice chapter.

\EH\textbf{Why the support memory alone supplies neither sign of RH.}
The map evaluating a globally supported function at $(0,Z_0)$ is
\[
 G_L(\mathcal O(\mathbb C^2))\longrightarrow G_L(\mathbb C),
 \qquad \widehat f\longmapsto\widehat{f(0,Z_0)}.
\]
It retains every lower support label. Applied to the original source,
$\widehat H(0,Z_0)=e$ if and only if $H_0(Z_0)=0$, equivalently
$\xi(1/2+iZ_0/2)=0$. Applying $\chi_L$ afterwards maps both supported
zero and every supported nonzero complex value to the same top label.
Thus that additional map cannot test which $Z_0$ is a zero. This is
proved by its displayed formula, without discarding the retained
support boundary.

There is also an explicit test of the amount of information in
\eqref{eq:eh-moments}'s support. Take the positive measure
$\delta_1$ on $[0,\infty)$ in the same cosine--heat integral. It gives
$K_t(Z)=e^t\cos Z$, all of whose even coefficients have exactly that
same nonzero sign pattern for all $m,n$. Every zero is real for every
real $t$, since $\cos Z=0$ is equivalent to $e^{2iZ}=-1$ and hence
$\Im Z=0$. Alternatively $Q_t(Z)=Z^2-2t$ gives the preceding
nontrivial collision and monodromy at a real-root time-zero fibre.
These test families are connected to the actual family by the common
linear integral operator $\mu\mapsto\int e^{tu^2}\cos(Zu)\,d\mu$;
their measures and values are stated, not substituted for $\Phi(u)du$.
They show that coefficient support, local invertibility at $e$, and
collision monodromy, without the arithmetic moment values, do not
force a failure of real-rootedness at zero.

\EH\textbf{The arithmetic threshold and the exact outstanding calculation.}
For the original $H_t$, the de Bruijn--Newman theorem gives real zeros
exactly for $t\ge\Lambda$. Rodgers--Tao prove $\Lambda\ge0$
(Theorem1.1; original TeX label \texttt{main}). Platt--Trudgian's
\href{https://arxiv.org/abs/2004.09765}{The Riemann hypothesis is true up
to $3\cdot10^{12}$}, section on the de Bruijn--Newman constant, original
TeX lines135--149, proves the corollary $\Lambda\le0.2$ using
\href{https://arxiv.org/abs/1904.12438}{D.H.J. Polymath}, Table1.
These cited results give $\mathrm{RH}\iff\Lambda=0$.
The calculations above produce neither a nonreal zero of the actual
$H_0$ nor a negative value of its whole Weil form. They identify the
exact retained endpoint and support memory and the arithmetic
evaluation map through which a proposed counterexample must pass.
The unital-localization receiver is $L$, and the precise local collision
above has a real double root.

\begin{figure}[p]
\centering
\begin{tikzpicture}[>=Stealth,node distance=1.1cm,
 box/.style={draw,rounded corners,align=center,text width=11.9cm,inner sep=7pt}]
\node[box] (a) {$\tau_{\rm f}<1/8$, $z=\sqrt{1-8\tau_{\rm f}}$\\
$\gamma=(z^{-1},-3z/2,13z^2/2)$, $F\gamma=(-z^2/4,0,0)$};
\node[box,below=of a] (b) {$r=-z/2$, $Z=-2ir=iz$, $h=4\tau_{\rm f}-1/2$\\
$Q_h(Z)=Z^2-2h=2P_{a,0,0}(r)$};
\node[box,below=of b] (c) {$\tau_{\rm f}\uparrow1/8$ maps to $h\uparrow0$\\
The affine coordinate $x$ escapes; both heat roots reach $Z=0$.};
\node[box,below=of c] (d) {$\mathcal A_0=\mathbb C[\eta]/(\eta^2)$\\
Loop around $h=0$: $M_{\rm loop}=\left(\begin{smallmatrix}0&1\\1&0\end{smallmatrix}\right)$,
 $M_{\rm loop}^*M_{\rm loop}=I$};
\node[box,below=of d] (e) {Split chart boundary:
$G_L(\mathbb C[z,z^{-1}])\amalg_{G_L(\mathbb C[z])}G_L(\mathbb C)=L$\\
All support labels remain. Complex amplitudes do not survive this receiver.};
\draw[->] (a)--(b);\draw[->] (b)--(c);\draw[->] (c)--(d);\draw[->] (d)--(e);
\end{tikzpicture}
\caption{Exact endpoint maps of EH2--EH7. The final arrow means applying
the split localization construction to the retained chart, not an
identification of the dual-number algebra with $L$. The labels give the
original coefficients and the complete time conversion. The heat
coordinate convention is Rodgers--Tao's; the polynomial orbit is the
retained programme calculation rederived here.}
\end{figure}
